\chapter{Measurement Updates, Gating and the Pose Sensor}
\label{ch:updates}
\chaptermark{Measurement Updates} % short running header

\begin{objectives}
\begin{itemize}
  \item Apply the generic manifold update to concrete sensors, deriving $H$ by the elementary Jacobians.
  \item Understand which state directions a measurement observes, from the nullspace of $H$.
  \item Implement the mock pose sensor of Milestone 6 and derive its Jacobian.
  \item Gate measurements with the $\chi^2$ test, and understand how gating fails in an overconfident filter.
  \item Understand how a pose sensor makes velocity and biases observable.
\end{itemize}
\end{objectives}

% ------------------------------------------------------------------------------
\section{The generic update}

\begin{convention}[measurement update, AGENTS.md F3]
\begin{enumerate}
  \item Residual $r = z\bminus h(\hat x)$: plain subtraction for vector measurements; for the pose measurement
        $r = [\,p_{meas} - \hat p;\ \Log_q(\hat q\inv\qprod q_{meas})\,]$.
  \item $H = \partial r/\partial\delta x$ at $\delta x = 0$ using C4 (\cref{sec:manifold-update}), derived by hand and verified by
        finite differences.
  \item $S = HPH\T + R_z$; $K = PH\T S\inv$ by a Cholesky/LDL$^\mathsf{T}$ solve, never an explicit inverse.
  \item Mahalanobis gate $r\T S\inv r < \chi^2_{0.95}(\dim r)$ from a \code{constexpr} table. Rejections are counted and
        logged, never silently dropped.
  \item Joseph form $P \gets (I - KH)P(I - KH)\T + KR_zK\T$, then injection and reset (\cref{sec:reset}).
\end{enumerate}
\end{convention}

Everything specific to a sensor is in $h$ and $H$. The rest of this section derives them for two sensors that the
project does not use, to demonstrate the technique, before turning to the pose sensor it does use.

% ------------------------------------------------------------------------------
\section{Worked example: position with a lever arm}

A GNSS antenna mounted at a known offset $\ell_b$ from the IMU, expressed in the body frame, measures the
antenna position:
\begin{equation}
  z = h(x) + n = p_{wb} + R_{wb}\,\ell_b + n, \qquad n \sim \N(0, \sigma^2I_3).
\end{equation}
The residual is a plain difference, $r = z - (\hat p + \hat R\ell_b)$. To find $H$, perturb the nominal state with
$\bplus$ and expand:
\begin{equation*}
  h(\hat x\bplus\delta x) = \hat p + \delta p + \hat R\Exp(\delta\theta)\ell_b
  \approx \hat p + \delta p + \hat R\ell_b + \hat R\sx{\delta\theta}\ell_b
  = h(\hat x) + \delta p - \hat R\sx{\ell_b}\delta\theta .
\end{equation*}
Reading off the coefficients (the rotation-action row of \cref{tab:jacobians}),
\begin{equation}
  H = \begin{bmatrix} I_3 & 0_3 & -\hat R\sx{\ell_b} & 0_3 & 0_3 \end{bmatrix}.
\end{equation}
The attitude column is non-zero because the antenna swings around the IMU when the body rotates. A
lever-arm position measurement therefore carries attitude information, but only about rotations that move
the antenna: the nullspace of $-\hat R\sx{\ell_b}$ is the direction $\ell_b$, since rotating about the lever arm itself
does not move the antenna.

% ------------------------------------------------------------------------------
\section{Worked example: a direction sensor}

A magnetometer, or an accelerometer used as an inclinometer at rest, measures a known world vector $m_w$ in
body coordinates:
\begin{equation}
  z = R_{wb}\T m_w + n .
\end{equation}
With the inverse-action row of \cref{tab:jacobians}, $h(\hat x\bplus\delta x) = \Exp(-\delta\theta)\hat R\T m_w \approx \hat R\T m_w +
\sx{\hat R\T m_w}\delta\theta$, so
\begin{equation}
  H = \begin{bmatrix} 0_3 & 0_3 & \sx{\hat R\T m_w} & 0_3 & 0_3 \end{bmatrix}.
\end{equation}
$\sx{v}$ has rank two and nullspace $\mathrm{span}(v)$, so a direction measurement observes two of the three attitude
degrees of freedom. The rotation about the measured direction is unobservable. For gravity this is the
familiar statement that an accelerometer at rest gives roll and pitch but not yaw.

\begin{keyidea}[the nullspace of $H$ tells you what a sensor cannot see]
Any error direction $\delta x \in \nullsp(H)$ produces no change in the predicted measurement, so the update cannot
correct it. Reading the nullspace of each new $H$ is a quick sanity check of both the derivation and your
intuition about the sensor.
\end{keyidea}

% ------------------------------------------------------------------------------
\section{The pose sensor of Milestone 6}

Before introducing a camera, the project tests the whole ESKF with a ``mock sensor'' that measures the pose
directly. Ground-truth poses, decimated to about 10~Hz, are corrupted with synthetic noise of known statistics.

\begin{convention}[pose measurement, AGENTS.md C6]
\begin{equation}
  p_{meas} = p_{wb} + n_p, \qquad q_{meas} = q_{wb}\qprod\Exp_q(n_\theta), \qquad
  R_z = \diag(\sigma_p^2 I_3,\ \sigma_\theta^2 I_3).
\end{equation}
\end{convention}

The attitude noise is applied on the right, in the body frame, which is exactly the local perturbation. That
choice makes the residual \eqref{eq:manifold-residual} statistically clean: if the nominal state were exact, the
attitude residual would be exactly $n_\theta$. To generate the noise, draw $n_p$ and $n_\theta$ from seeded
\code{std::mt19937\_64} generators (AGENTS.md E5) with the same $\sigma_p$, $\sigma_\theta$ that appear in $R_z$.

\begin{derivex}{the pose measurement Jacobian $H \in \R^{6\times15}$}{sol:Hpose}
Using \eqref{eq:manifold-H}, derive $H$ for $h(x) = (p_{wb}, q_{wb})$ with the residual of AGENTS.md F3.
\begin{enumerate}
  \item The position rows are a plain difference. Which state blocks does $\hat p + \delta p$ depend on?
  \item For the attitude rows, evaluate $\Log_q\big(\hat q\inv\qprod(\hat q\qprod\Exp_q(\delta\theta))\big)$ and simplify.
        Which elementary Jacobian applies if you instead linearize around a non-zero residual?
  \item Write $H$ in the frozen column order and state its nullspace. Which states are \emph{not} directly measured,
        and how can they still be corrected?
\end{enumerate}
\textbf{Check:} the finite-difference test of $r(\delta x) = h(\hat x\bplus\delta x)\bminus h(\hat x)$ at a non-trivial state, as
required by Milestone 6.
\end{derivex}

\paragraph{How unmeasured states get corrected.} The pose sensor measures neither velocity nor biases, yet the
update corrects both. The mechanism is the cross-covariance built up by propagation. Between updates, a velocity
error produces a growing position error, so $P_{pv}$ becomes large; an accelerometer bias produces a velocity error
and then a position error, so $P_{p b_a}$ grows; a gyro bias produces an attitude error, so $P_{\theta b_g}$ grows. The gain
$K = PH\T S\inv$ has non-zero rows for every state correlated with the measured ones, so each pose update corrects
them in proportion to those correlations.

\paragraph{Why the biases converge.} Observability here has an intuitive reading. The gyro bias is the slope of the
difference between the integrated gyro and the measured attitude, which a sequence of attitude measurements reveals.
The velocity is the slope of the measured position. The accelerometer bias is the difference between the
acceleration implied by successive velocities and the specific force rotated with the (measured) attitude, so it
too becomes observable. Even at rest it is separable from tilt, because the attitude is measured directly. With
pose measurements every state of the IMU error state is observable, and their covariances converge.

% ------------------------------------------------------------------------------
\section{Gating}
\label{sec:gating}

Real measurements include outliers: a misassociated feature, a reflection, a moving object. A single outlier with a
large residual can drag the estimate far away and corrupt the covariance. The $\chi^2$ test of \cref{prop:innovation}
rejects measurements that are statistically implausible given the filter's own uncertainty:
\begin{equation}
  \gamma = r\T S\inv r \ \stackrel{?}{<}\ \chi^2_{1-\alpha}(\dim r).
\end{equation}

\begin{table}[h]
\centering\small
\begin{tabular}{@{}c c c c c c c@{}}
\toprule
$\dim r$ & 1 & 2 & 3 & 4 & 5 & 6 \\
\midrule
$\chi^2_{0.95}$ & 3.841 & 5.991 & 7.815 & 9.488 & 11.070 & 12.592 \\
$\chi^2_{0.99}$ & 6.635 & 9.210 & 11.345 & 13.277 & 15.086 & 16.812 \\
\bottomrule
\end{tabular}
\caption{Gate thresholds. A longer table, up to 600 degrees of freedom, is in \cref{app:tables}.}
\label{tab:chi2-short}
\end{table}

With a 95\% gate, a consistent filter rejects 5\% of perfectly good measurements. That is the price of rejecting
most outliers, and it is why rejections are counted and logged: a rejection rate close to 5\% is a sign of
health, a rate near zero suggests an underconfident filter (or a gate that is too loose to be useful), and a
rate far above 5\% indicates overconfidence or a model error.

\begin{pitfall}[the death spiral]
Gating assumes the filter is consistent. If $P$ is too small, good measurements produce large $\gamma$ and are
rejected; with no corrections $P$ should grow, but if the model's process noise is also too small it grows slowly
while the true error grows quickly, and more measurements are rejected. The filter then dead-reckons with false
confidence until it diverges. When the rejection rate jumps, suspect overconfidence before suspecting outliers.
\end{pitfall}

The threshold depends on the dimension of $r$. A pose measurement uses $\chi^2_{0.95}(6)$, a single image
observation $\chi^2_{0.95}(2)$, and a projected MSCKF feature residual $\chi^2_{0.95}(2M-3)$, where $M$ is the number of
observations (\cref{ch:msckfupdate}). Because the quantiles grow with the dimension, the table is generated offline
and stored as a \code{constexpr} array, which avoids depending on a statistics library in the core.

% ------------------------------------------------------------------------------
\section{Pseudo-measurements and timing}

\paragraph{Pseudo-measurements.} Knowledge that is not a sensor reading can be fed to the filter as if it were one. The most useful is the \term{zero-velocity
update} (ZUPT): when the vehicle is known to be at rest (detected, for example, from a very small variance of the accelerometer and gyroscope readings over
a window), apply the measurement $z = 0 = v_{wb} + n$ with $H = [0\ \ I_3\ \ 0\ \ 0\ \ 0]$ and a small noise covariance. It stops velocity drift, and through the
correlations it also refines tilt and biases. A ZUPT applied when the vehicle is in fact moving is a large, confident, wrong measurement; the detection must
be conservative. Similar pseudo-measurements express a known constant altitude, a non-holonomic constraint (a car does not slide sideways), or a
known landmark.

\paragraph{Measurement timing.} A measurement taken at time $t_z$ must be applied to the state at $t_z$, after propagating exactly to that time
(\cref{ch:strapdown}). Measurements that arrive \emph{late}, after the filter has already propagated past their timestamp, are a problem for every filter:
either the filter buffers the IMU data and re-propagates from $t_z$, or it uses a stochastic clone of the state at $t_z$ (\cref{ch:cloning}), which is exactly the
MSCKF's approach for camera measurements. On EuRoC, data are processed offline in timestamp order, so late measurements do not occur.

% ------------------------------------------------------------------------------
\section{Testing the update}

\begin{projectlink}[Milestone 6]
The exit criteria test properties that hold for any correct implementation:
\begin{itemize}
  \item $H$ passes the finite-difference test at a non-trivial state;
  \item $P_{prior} - P_{post}$ is positive semi-definite (an update never adds uncertainty);
  \item after injection and reset all quaternions are unit, and a second update with the same $z$ produces a residual
        consistent with the first correction;
  \item as $R_z \to 0$ the posterior pose converges to $z$;
  \item with constant synthetic biases, the bias estimates converge to within their own $3\sigma$;
  \item update, injection and reset do not allocate;
  \item a full MH\_01 run with noisy ground-truth updates neither diverges nor produces NaN, and the $\chi^2$ rejection
        count is logged.
\end{itemize}
\end{projectlink}

The ``second update'' test is subtle and worth understanding. After the first update, the nominal pose has moved
toward $z$ by $K r$, so the second residual should be approximately $(I - HK)r$, which is smaller. If it is instead
larger, or has the opposite sign, the injection used the wrong side of the quaternion product or the wrong sign of
$H$.

\begin{takeaways}
\begin{itemize}
  \item Derive $H$ by perturbing $h(\hat x\bplus\delta x)$, expanding to first order and reading off coefficients.
  \item The nullspace of $H$ lists the error directions a sensor cannot see.
  \item The pose sensor's attitude noise is a local perturbation, matching the residual $\Log_q(\hat q\inv\qprod q_{meas})$.
  \item Unmeasured states are corrected through cross-covariances built up by propagation.
  \item Gate with $\chi^2_{0.95}(\dim r)$, log the rejections, and treat a high rejection rate as a symptom of overconfidence.
\end{itemize}
\end{takeaways}

\begin{exercises}
  \item For the lever-arm example, compute $S$ and $K$ for a state with
        \[ P = \diag(\sigma_p^2I,\ \sigma_v^2I,\ \sigma_\theta^2I,\ \sigma_{b_a}^2I,\ \sigma_{b_g}^2I). \]
        Which components of $\delta\theta$ does the update correct?
  \item Derive $H$ for a range measurement to a known beacon $b_w$: $z = \|p_{wb} - b_w\| + n$. What is its nullspace?
  \item Simulate the death spiral: run a 1D constant-velocity filter with $Q$ underestimated by a factor of 100 and a
        95\% gate. Plot the rejection rate and the error over time.
  \item Show that with a pose measurement of a static vehicle, the horizontal accelerometer bias is observable.
        (Hint: what does the attitude measurement remove from the tilt/bias ambiguity of \cref{ch:strapdown}?)
  \item Prove that $P_{prior} - P_{post} = KSK\T \succeq 0$ for the optimal gain, and explain why this test can fail in floating
        point if the simple covariance update is used.
\end{exercises}
% Solutions to the exercises of Chapter 18.
\begin{solutions}
  \item With $H = [\,I_3\ \ 0\ \ {-\hat R\sx{\ell_b}}\ \ 0\ \ 0\,]$ and the block-diagonal $P$, only the position and attitude blocks contribute:
        \[
          S = \sigma_p^2I + \sigma_\theta^2\hat R\sx{\ell_b}\sx{\ell_b}\T\hat R\T + \sigma^2I = \alpha\,I - \sigma_\theta^2\,(\hat R\ell_b)(\hat R\ell_b)\T,
          \qquad \alpha = \sigma_p^2 + \sigma^2 + \sigma_\theta^2\|\ell_b\|^2,
        \]
        using $\sx{\ell}\sx{\ell}\T = -\sx{\ell}^2 = \|\ell\|^2I - \ell\ell\T$. The gain blocks are $K_p = \sigma_p^2S\inv$, $K_\theta = \sigma_\theta^2(-\hat R\sx{\ell_b})\T S\inv = \sigma_\theta^2\sx{\ell_b}\hat R\T S\inv$,
        and zero for velocity and biases. Since $S\inv(\hat R\ell_b) \propto \hat R\ell_b$ and $\sx{\ell_b}\hat R\T\hat R\ell_b = 0$, the attitude gain simplifies to
        \[
          K_\theta = \frac{\sigma_\theta^2}{\sigma_p^2 + \sigma^2 + \sigma_\theta^2\|\ell_b\|^2}\,\sx{\ell_b}\hat R\T .
        \]
        Its range is the range of $\sx{\ell_b}$: the update corrects only the attitude components perpendicular to the lever arm, never the rotation
        about $\ell_b$ itself, which does not move the antenna. Velocity and biases are not corrected here because this prior has no cross-covariances;
        after some propagation they would be.

  \item With $h(x) = \|p - b_w\|$ and $\hat\rho = \|\hat p - b_w\|$, perturbing $p = \hat p + \delta p$ gives $h \approx \hat\rho + u\T\delta p$ with $u = (\hat p - b_w)/\hat\rho$, so
        \[
          H = [\,u\T\ \ 0\ \ 0\ \ 0\ \ 0\,] \in \R^{1\times15} .
        \]
        Its nullspace is 14-dimensional: the two position directions perpendicular to the line of sight (moving on the sphere around the beacon), and
        all velocity, attitude and bias errors. One range observes one number; position perpendicular to the line of sight becomes observable only
        through several beacons or through motion and the correlations it builds.

  \item Simulate a constant-velocity target (acceleration density $\sqrt{0.5}$, position measurements at 10~Hz with unit variance) and run the
        filter with $Q$ divided by 100 and a $\chi^2_{0.95}(1)$ gate. Typical results, per 30-second window:
        \begin{center}\small
        \begin{tabular}{@{}l l l@{}}
        \toprule
        filter & rejection rate & RMS position error \\
        \midrule
        correct $Q$, gated & 2--8\% & 0.4--0.65~m \\
        $Q/100$, gated & 79\%, then 100\% & 26, 126, 243, \dots, 3246~m \\
        $Q/100$, not gated & 0\% & 0.8--1.7~m \\
        \bottomrule
        \end{tabular}
        \end{center}
        The overconfident filter predicts innovations that are far too small, so good measurements fail the gate; each rejection removes the
        correction the filter needed, the error grows, more measurements are rejected, and within half a minute the filter ignores all data and
        dead-reckons. Its reported $\sigma$ (about 200~m at the end) stays far below its actual error. Without the gate the same mistuned filter is
        merely suboptimal: its error stays bounded. A rejection rate far above the design value is the alarm signal (\cref{sec:gating}).

  \item At rest, the accelerometer measures $\hat R\T(-g_w) + b_a$, and to first order a tilt error and a horizontal bias change this reading in the same way:
        in the velocity error dynamics both enter through the horizontal specific-force error, so from accelerations alone only their combination is
        observable (\cref{ch:strapdown}). With a pose sensor, the position measurements make the velocity error and its rate observable (two
        differentiations, as in \cref{sec:observability-linear}), while the attitude measurement observes $\delta\theta$ directly. The rate of the velocity
        error is (tilt term) $+$ (bias term); with the tilt term known from the attitude measurement, the bias term, and hence the horizontal bias, is
        determined. The vertical bias was observable already (it changes the measured magnitude of gravity, which the position measurements expose),
        and the gyro bias is observable from the drift of the attitude between measurements. With a pose sensor, every IMU state is observable,
        which is why the covariances converge in Milestone 6.

  \item For the optimal gain, $P_{\mathrm{prior}} - P_{\mathrm{post}} = KHP = KSK\T$ (\cref{thm:kf-update}), and $KSK\T \succeq 0$ because it is a congruence of $S \succ 0$
        (\cref{prop:psd}). In floating point the simple update $P_{\mathrm{post}} = (I - KH)P$ is neither symmetric nor guaranteed PSD (Exercise 12.3), and
        $P_{\mathrm{prior}} - P_{\mathrm{post}}$ then has rounding-level asymmetry and can have slightly negative eigenvalues, notably in directions the
        measurement does not observe, where the true difference is exactly zero and the computed one is $\pm u\|P\|$. A test must therefore symmetrize the
        difference and allow a tolerance: $\lambda_{\min}\big(\tfrac12(D + D\T)\big) \ge -c\,u\,\|P\|$ with a small constant $c$. Using the Joseph form and
        symmetrizing after each update makes $P_{\mathrm{post}}$ itself PSD, but the difference still needs the tolerance.
\end{solutions}

